% Complete Section 7.7, original PDF 101-103 / printed 91-93.
\begingroup\emergencystretch=2em
% block: cubical-elements-heading page:101 type:section
\subsection{{Cubical elements}}\label{sec:cubical-elements}

% block: cubical-elements-introduction pages:101,102 type:ordinary-paragraph
{Although we have only discussed simplicial triangulations to this point, finite element spaces based on meshes of $n$-dimensional boxes (products of intervals) are also popular. It turns out that there are two families of finite element spaces on boxes which are in many ways parallel to the $\mathcal P_r\Lambda^k$ and $\mathcal P_r^-\Lambda^k$ spaces on simplices and which we summarize in this section.}

% block: tensor-product-form-family page:102 type:ordinary-paragraph
{The simpler of the two families is denoted by $\mathcal Q_r^-\Lambda^k$. It is defined on an arbitrary mesh of boxes $\mathcal T_h$ in $\mathbb R^n$ and for every $r>0$ and $0\leq k\leq n$. To define the space, we need to supply shape functions and DOFs on an arbitrary box $T\subset\mathbb R^n$. This can be done using the appropriate notion of a tensor product of chain complexes. Here we quickly recall the results, which are extracted from \BookCite{9}, to which we refer for details.}

% block: tensor-product-shape-functions page:102 type:ordinary-paragraph
{Without loss of generality, we take $T=I^n$, where $I=[0,1]$. Then the shape function space is}
\[
\mathcal Q_r^-\Lambda^k(I^n)=\bigoplus_{1\leq\sigma_1<\cdots<\sigma_k\leq n}\left[\bigotimes_{i=1}^n\mathcal P_{r-\delta_{i,\sigma}}(I)\right]dx^{\sigma_1}\wedge\cdots\wedge dx^{\sigma_k},
\]
{where}
\[
\delta_{i,\sigma}=\begin{cases}1,&i\in\{\sigma_1,\ldots,\sigma_k\},\\0&\text{otherwise}.\end{cases}
\]
{In the case $k=0$, this space is understood to be}
\[
\mathcal Q_r^-\Lambda^0(I^n)=\bigotimes_{i=1}^n\mathcal P_r(I),
\]
{i.e., the space conventionally referred to as $\mathcal Q_r(I^n)$. This definition makes sense also if $r=0$, so $\mathcal Q_0^-\Lambda^0(I^n)=\mathcal Q_0(I^n)=\mathbb R$ is the space of constant functions. For $k>0$, $\mathcal Q_0^-\Lambda^k(I^n)=0$. We also allow $n=0$ in this definition, i.e., when $I^n$ reduces to a single point. We then understand $\mathcal Q_r^-\Lambda^0(I^0)$ to be the space $\mathbb R$ of constants. It is easy to count the dimension of these spaces:}
\[
\dim\mathcal Q_r^-\Lambda^k(I^n)=\binom nk(r+1)^{n-k}r^k,\qquad0\leq k\leq n,\ r\geq0.
\]

% block: tensor-product-degrees-of-freedom page:102 type:ordinary-paragraph
{Finally we specify a set of DOFs for $\mathcal Q_r^-\Lambda^k(I^n)$ ($r\geq1$, $0\leq k\leq n$),}
\begin{equation}\label{eq:tensor-product-form-dofs}
v\mapsto\int_f\operatorname{tr}_fv(x)\wedge q(x),\qquad q\in\mathcal Q_{r-1}^-\Lambda^{d-k}(f),
\end{equation}
{for each face $f$ of $I^n$ of degree $d\geq k$. As shown in \BookCite{9} their unisolvence follows from tensor product considerations. For each $r\geq1$, the resulting finite element spaces combine to give a finite element de Rham subcomplex:}
\[
0\to\mathcal Q_r^-\Lambda^0(\mathcal T_h)\xrightarrow{d}\mathcal Q_r^-\Lambda^1(\mathcal T_h)\xrightarrow{d}\cdots\xrightarrow{d}\mathcal Q_r^-\Lambda^n(\mathcal T_h)\to0.
\]
{Note that, like the $\mathcal P_r\Lambda^k$ sequence, the polynomial degree remains constant across the complex. Moreover, the canonical projections associated to the DOFs form a cochain map.}

% block: serendipity-shape-space pages:102,103 type:ordinary-paragraph
{The second family of cubical finite elements was derived in \BookCite{7}, using the tools of FEEC. We refer to it as the \emph{serendipity family} of finite element differential forms and denote the finite element space by $\mathcal S_r\Lambda^k$. It has a somewhat complicated definition of the shape functions. First we define the space of polynomial differential forms $\mathcal H_{r,s}\Lambda^k$ as the span of those monomials $(x^1)^{\alpha_1}\cdots(x^n)^{\alpha_n}\,dx^{\sigma_1}\wedge\cdots\wedge dx^{\sigma_k}$ which are}
\begin{enumerate}
\item{of polynomial degree $r$ (i.e., $\alpha_1+\cdots+\alpha_n=r$),}
\item{are linear in at least $s$ variables, not counting the variables $x^{\sigma_i}$, i.e.,}
\[
\#(\{\alpha_i=1\}\setminus\{\sigma_1,\ldots,\sigma_k\})\geq s.
\]
\end{enumerate}
{Then we take as shape function space the space of polynomial differential forms}
\[
\mathcal S_r\Lambda^k:=\mathcal P_r\Lambda^k\oplus\bigoplus_{s\geq1}[\kappa\mathcal H_{r+s-1,s}\Lambda^{k+1}\oplus d\kappa\mathcal H_{r+s,s}\Lambda^k].
\]

% block: serendipity-degrees-of-freedom page:103 type:ordinary-paragraph
{In contrast to this intricate definition of the shape functions, the DOFs for the $\mathcal S_r\Lambda^k$ family are mercifully simple. To a face $f$ of dimension $d$ we associate the DOFs}
\[
u\mapsto\int_f(\operatorname{tr}_f\omega)\wedge\mu,\qquad q\in\mathcal P_{r-2d}\Lambda^{d-k}(f),
\]
{where the weighting functions are taken from ordinary full polynomial spaces, but their degrees drop by two each time the dimension of the face increases by one. In \BookCite{7}, we show that these DOFs are unisolvent, and the main properties of the family are determined. In particular, it is shown that the assembled finite element spaces combine into a de Rham subcomplex with commuting canonical projections:}
\[
0\to\mathcal S_r\Lambda^0(\mathcal T_h)\xrightarrow{d}\mathcal S_{r-1}\Lambda^1(\mathcal T_h)\xrightarrow{d}\cdots\xrightarrow{d}\mathcal S_{r-n}\Lambda^n(\mathcal T_h)\to0.
\]
{Note that the polynomial degrees decrease as for the $\mathcal P_r\Lambda^k$ family.}

% block: further-cubical-families-geometric-distortion page:103 type:ordinary-paragraph
{We remark that some additional families of cubical finite element differential forms have been proposed at least in two or three dimensions, sharing many of the same properties. In particular, we note the TNT element of Cockburn and Qiu \BookCite{33} and the trimmed serendipity family of Gillette and Kloefkorn \BookCite{44}. Another remark is that for all these cubical elements, including these discussed above, there is a major loss of accuracy if the cube is mapped not to a box or a parallelepided but to a more general shape (such as a convex quadrilateral in two dimensions, or some sort of hexahedron in three dimensions). This effect, which becomes more pronounced the larger the form degree of the elements, was studied in \BookCite{9}. Some elements which avoid this have recently been constructed by Arbogast and collaborators \BookCite{3}, \BookCite{4}.}
